\section{Crossings, cylinder barriers, and positive nesting}\label{c6:geometry}

The strip estimate supplies the mass of a crossing with one endpoint fixed. We first turn it into boundary escape and confined-crossing estimates. Averaging the cuts of several such crossings then produces essential polygons of positive mass on a cylinder. These polygons give height barriers; a separate cut-link identity transfers positive polygon nesting to the number of links used by free-boundary paths.

We use ordinary honeycomb paths at weight \(x^{\#\text{ triangle centers visited}}\), usually \(x=\kappa\), as well as unoriented unrooted simple polygons with the same weight. Put \(d_0=\sqrt3/2\). Take a horizontal side direction of the triangular lattice, \(n_0\) vertically up and \(n_j\) obtained from \(n_0\) by clockwise rotations through \(j\pi/3\). Lines normal to \(n_j\) which bound triangular rows have level spacing \(d_0\); crossing them is through midpoints of edges of the tiling in these lines (a visit at a midpoint is transverse). A strip of height \(k d_0\) normal to any \(n_j\) has the boundary-to-boundary mass \(B_k\) from each fixed midpoint on one rim. Thus the strip estimate above gives \(B_k\asymp(k+1)^{-1/4}\) (we put \(B_0=1\) for zero bridge gaps).

\paragraph{Criticality with the boundary-port convention.}
We adapt the Hammersley--Welsh bridge decomposition, as used in \cite[Section~3]{DuminilCopinSmirnov2012}, to the present strip and endpoint convention. In particular \(\kappa\) really is the critical single-walk weight (radius of convergence of the usual length series from a fixed honeycomb vertex). Here are details to account for the change of boundary conventions. At \(x=\kappa\), summing \(B_k\) over heights diverges, with starting triangle fixed, hence also the unrestricted susceptibility (each walk with its initial and final triangles specified incurs at most bounded multiplicity). At \(x<\kappa\), measure heights of triangle centers in units \(d_0/3\). A path starting from a fixed center and making a weak bridge (staying inclusively between the heights of its initial and final vertices, whose difference has absolute value \(k>0\)) has total mass at most \(C_x e^{-c_x k}\), using here \(x^{\#\text{steps}}\). Indeed for the increasing direction one can attach a bottom end on the bottom line of the initial vertex's row, and a top end on the top line of the last vertex's row. This uses no extra triangles if the triangle already has a ray in that boundary direction, and otherwise uses just an adjacent triangle strictly below or strictly above in its row, with a deterministic choice. The resulting walk is simple and a strip bridge, with known start and bounded multiplicity; thus the mass at \(\kappa\) was bounded, and there is exponential suppression by the minimal length. The decreasing case is reflected. A weak upward half-space walk from its lowest endpoint splits into such bridges by taking its last highest vertex, then in the remainder its last lowest vertex, and so on until exhausted. Their positive spans are decreasing strictly except possibly equality of the first two (consecutive heights along any single step differ). The total half-space mass is therefore finite by summability of the bridge bounds and \(\prod_{k\ge1}(1+C_xe^{-c_xk})^2<\infty\). Any unrestricted walk splits at a lowest vertex into two weak half-space walks rooted there; one may translate that vertex to fixed representatives of the two types, recovering the translation by the prescribed original start. This proves convergence. The two triangle types used here are equivalent by symmetry.

\subsection{Flat and turned crossings}\label{c6:subheading-11}

\begin{proposition}[Boundary crossings]\label{prop:crossings}
For a fixed boundary port of a convex pure-sided tiled domain, including a straight strip or half-plane by exhaustion, the critical mass of completed paths with diameter at least $r$ is at most $C(1+r)^{-1/4}$. In a pure half-plane the mass $K_l$ between two boundary ports at lattice distance $l\ge1$ is decreasing in $l$ and satisfies $K_l\asymp l^{-5/4}$. All constants are independent of the domain and the chosen port.
\end{proposition}
\begin{proof}
We first derive escape estimates from the contour identity, then obtain the two-sided half-plane bound.

\paragraph{Positive contour comparison.}
The local cancellation in Section~\ref{sec:strip-flux} applies in every finite convex tiled polygon. A path from a boundary-edge midpoint has phase \(e^{i\beta W}\), where \(\beta=3/8\) and \(W\) is its total turn. Convexity bounds the completed turn by \(|W|\le\pi\), so the signed exit identity gives
\[
 \sum_{\gamma:\,p\to\partial D}\kappa^{\ell(\gamma)}
       \cos(\beta W(\gamma))=1,
 \qquad c\le\cos(\beta W(\gamma))\le1,
 \qquad c=\cos(\beta\pi)>0.
\] Thus cutting down to a convex subdomain sharing the source with unchanged initial direction, the mass of complete paths lost (those not themselves paths to a shared exit inside the subdomain) is at most \(c^{-1}\) times the mass from the source to the new sides in the subdomain, by subtraction. Such upper bounds apply by exhaustion if the large domain is unbounded.

\paragraph{Triangle escape.}
Let
\[
 T(p,n_j,t)=\{n_j\cdot(z-p)\ge0,\quad n_{j\pm1}\cdot(z-p)\le t\},
 \qquad t/d_0\in\mathbb Z_{\ge0}+1/2
\]
be a centered triangle based at a midpoint \(p\) (of a tiling edge normal to \(n_j\), i.e., in the base line with inward normal \(n_j\)). Exit mass from \(p\) to each sloping side is the same, say \(S(t)\). If \(A_T\) is the base-return mass, \(2\cos(\beta\pi/3)S=1-cA_T\); thus \(S\) decreases. We have
\begin{equation}\label{eq:triangle-escape}
 S(t)\asymp(1+t)^{-1/4}.
\end{equation}
Indeed embed the triangle in a strip of height a larger multiple of order \(t\) to see \(B_k\lesssim S(t)\) by contour comparison. Conversely allow all midpoints \(p'\) in the central half of the base as source in that same triangle. Their exit mass to the sloping sides is at least \(2 S(t')\) where \(t'\) is a common fixed-factor enlargement of \(t\) (rounded) that would contain this triangle when centered at \(p'\), by the base-return identity. Reverse and sum: from each side endpoint the mass to those central sources is \(O(B_k)\) for a strip width \(k d_0\asymp t\), by cutting the triangle parallel to that side strictly before all central targets. There are \(O(t)\) side endpoints and order \(t\) sources, proving the other bound by monotonicity. Consequently in convex polygons (also straight strips or the half-plane), the mass from a boundary source of completed paths of diameter \(\ge r\) is \(O((1+r)^{-1/4})\): intersect with a centered \(T\) of small comparable size and bound new-side mass by its triangle sum. Also in a straight half-plane the base-return mass is exactly \(1/c\), by the exhaustion \(T\).

\paragraph{Monotonicity along a flat boundary.}
Write \(K_l\) for half-plane mass from a base midpoint to another at distance \(l>0\). It decreases with \(l\). For details take midpoints \(a<P<Q=P+1\) along the base of an upper half-plane, and put an artificial arc \(P Q\) below, a semicircle tangent to the normal rays. Sum the prefix fluxes from \(a\) through a large truncating triangle, taking just terms using the artificial arc once (with its turn included and modulus weight 1). At triangle centers the flux rule still holds: unfinished loops can be reversed, preserving the arc count; the source is outside any loop even with that arc, so turn increments are again opposite \(\pm4\pi/3\). Thus the entering flux supplied at the end of an arc traversal (before further continuation) equals in total the flux through outer exits of paths using the arc. No such path exits afterwards at either arc anchor (the incident vertex needed there has already been visited). Contributions ending on the distant sides vanish in the limit, bounded by products of an ordinary path mass to an anchor and the ordinary distant-side escape mass from the other anchor. Put \(z=e^{i\beta\pi}\). Entering flux at \(Q\) and \(P\) therefore gives
\[
 K_{P-a}+z^{-2}K_{Q-a}=z D_1+z^{-1}D_2,
\]
with \(D_1\ge0\) the mass using \(a\to P\), the arc, then \(Q\to b\) for \(b<a\); \(D_2\ge0\) consists of the others, namely that first traversal ending at \(b>Q\), or \(a\to Q\) then via \(P\) ending \(a<b<P\). These choices follow by non-crossing and adjacency, with each half-plane strand turning \(\pm\pi\) according to order. In these coordinates \(a\to P,a\to Q\) turn clockwise by total \(-\pi\), the arc from \(P\) to \(Q\) gives \(+\pi\) (opposite in reverse). Loops using the arc have their interior on the base line only between the two anchors, hence cannot contain the source. Taking imaginary parts after multiplication by \(z\) gives the precise sign comparison
\[
 (K_{P-a}-K_{Q-a})\sin(\beta\pi)
   =D_1\sin(2\beta\pi)\ge0.
\]
This proves monotonicity. The diameter bound applied to endpoints in a window of order \(l\), followed by monotonicity, gives \(K_l\lesssim l^{-5/4}\).

\paragraph{Averaged deep exits.}
We record the triangle estimate needed both here and in the later collar and sewing constructions. All averages over lattice levels below are uniform averages over the admissible levels in the indicated interval. For an exit from \(T(a,n_0,t)\) through the side with outward normal \(n_1\), let \(k(t,q)\) be its mass when the exit port \(p_1\) has depth \(q=n_0\cdot(p_1-a)\). For fixed \(0<\alpha<A\) and sufficiently large \(R\),
\begin{equation}\label{eq:triangle-marginal}
 \sum_{t\in[\alpha R,AR]}k(t,q)\le C_{\alpha,A}R^{-1/4}
\end{equation}
uniformly in \(q\). To prove this, translate along the base by lattice steps to align the exit ports and their side lines, then reverse the paths. The translated starts are distinct and lie on the base at distance at least \(\alpha R\) from the common side line. In the convex wedge between that line and the base, cut parallel to the side strictly before all those starts. Contour comparison bounds the resulting mass by a strip crossing of width comparable to \(R\), proving \eqref{eq:triangle-marginal}.

Since \(S(t)\asymp R^{-1/4}\) throughout this window, the mass of exits of depth below \(\delta R\), averaged over \(t\), is at most \(C_{\alpha,A}\delta R^{-1/4}\). Consequently, for sufficiently small fixed \(\delta>0\),
\begin{equation}\label{eq:deep-triangle-exit}
 \mathop{\mathrm{average}}_{t\in[\alpha R,AR]}
       \sum_{q\ge\delta R}k(t,q)\ge c_{\alpha,A}R^{-1/4}.
\end{equation}
The threshold and constants may depend on the fixed window. In particular the same statement applies to a small fixed window of scales relative to a later construction.

\paragraph{Three exits and the lower half-plane bound.}
Triangle concatenations and averaging over renewal cuts have precedents in \cite{GlazmanManolescu2017,KrachunPanagiotis2026}. Here the deep-exit restriction gives the unique-crossing order needed for a second moment on one fixed path measure.
Starting at \(a\), make a deep exit from \(T(a,n_0,t)\) at \(p_1\), then exit \(T(p_1,n_1,q)\) through its side of normal \(n_2\), at \(p_2\). Put \(X=n_1\cdot(p_2-p_1)\), and finally exit \(T(p_2,n_2,X)\) through the side of normal \(n_3=-n_0\). Both \(q,X\) are admissible half-odd sizes. The identities \(n_2=n_1-n_0\) give
\[
 n_2\cdot(p_2-a)=t,\qquad n_0\cdot(p_2-a)=X;
\]
thus the final side is on the original base. Figure~\ref{fig:geometry-three-exits} shows the three domains and the two intermediate lines.
\begin{figure}[htbp]
\centering
\begin{tikzpicture}[x=.98cm,y=.98cm,>=Latex,font=\small,
  region/.style={draw=black!55,thick},
  walk/.style={draw=blue!65!black,very thick},
  cut/.style={draw=black!60,densely dashed}]
  \coordinate (a) at (0,0);
  \coordinate (r) at (2.30940,0);
  \coordinate (p1) at (1.27017,1.8);
  \coordinate (p2) at (3.69504,2.4);
  \coordinate (b) at (6.3,0);

  \path[region,fill=blue!7] (-2.30940,0)--(r)--(0,4)--cycle;
  \path[region,fill=orange!13] (r)--(.23094,3.6)--(4.38786,3.6)--cycle;
  \path[region,fill=green!7] (r)--(7.85196,0)--(5.08068,4.8)--cycle;
  \draw[thick] (-2.65,0)--(8.2,0);
  \node[anchor=north] at (4.5,-.60) {original base: $n_0\cdot(z-a)=0$};

  \draw[cut] (2.54034,-.4)--(-.80829,5.4)
    node[above] {$n_1\cdot(z-a)=t$};
  \draw[cut] (2.07846,-.4)--(5.42709,5.4)
    node[above] {$n_2\cdot(z-a)=t$};

  \draw[<->,black!55,thin] (1.27017,.03)--(1.27017,1.77)
    node[midway,right=3pt] {$q$};
  \draw[<->,black!55,thin] (3.69504,.03)--(3.69504,2.37)
    node[midway,left=3pt] {$X$};

  \draw[walk,->] (a) .. controls (0,.65) and (-.18,1.02) .. (.35,1.35)
    .. controls (.68,1.52) and (.94,1.61) .. (p1);
  \draw[walk,->] (p1) .. controls (1.83,2.12) and (2.01,2.61) .. (2.52,2.58)
    .. controls (3.03,2.55) and (3.26,2.65) .. (p2);
  \draw[walk,->] (p2) .. controls (4.18,2.12) and (4.37,2.55) .. (4.9,2.44)
    .. controls (5.8,2.15) and (6.3,.8) .. (b);

  \node[align=center] at (-.2,1.9) {$T(a,n_0,t)$};
  \node[align=center] at (2.16,3.22) {$T(p_1,n_1,q)$};
  \node[align=center] at (5.08,2.9) {$T(p_2,n_2,X)$};
  \fill (a) circle (1.8pt) node[below left=2pt] {$a$};
  \fill (p1) circle (1.8pt) node[above left=2pt] {$p_1$};
  \fill (p2) circle (1.8pt) node[above left=3pt] {$p_2$};
  \fill (b) circle (1.8pt) node[below right=2pt] {$b$};
  \draw[->] (-2.35,3.65)--(-2.35,4.4) node[above] {$n_0$};
  \draw[->] (-2.35,3.65)--(-1.7,4.025) node[right] {$n_1$};
  \draw[->] (-2.35,3.65)--(-1.7,3.275) node[right] {$n_2$};
\end{tikzpicture}
\caption{The three exits producing a half-plane return from $a$ to $b$.
The depths are $q=n_0\cdot(p_1-a)$ and $X=n_0\cdot(p_2-a)$.
The first triangle lies before the first dashed line, the second lies
after that line and before the second, and the third lies after the
second line and above the original base. Thus their interiors are
disjoint, and each dashed line is crossed once by the concatenation.
The triangle boundaries have the indicated directions; the blue paths
are schematic and the lattice is omitted.}
\label{fig:geometry-three-exits}
\end{figure}


The interiors are disjoint. The second triangle satisfies
\(n_1\cdot(z-a)\ge t\), \(n_2\cdot(z-a)\le t\), while the third satisfies
\(n_2\cdot(z-a)\ge t\), \(n_0\cdot(z-a)\ge0\), and remains on or after the first line. Each of the lines \(n_1\cdot(z-a)=t\) and \(n_2\cdot(z-a)=t\) is crossed exactly once. Hence the concatenation is simple and recoverable for a fixed \(t\).

Average first over \(t\in[R,2R]\). Equation~\eqref{eq:deep-triangle-exit} retains first-exit mass \(\gtrsim R^{-1/4}\) at \(q\ge\delta R\). The second triangle then gives total two-exit mass \(\gtrsim R^{-1/2}\). The part with \(X<\eta R\) is at most
\[
 O(R^{-5/4})\sum_{\delta R\le q\le4R}\sum_{X<\eta R}k(q,X)
       \lesssim_\delta\eta R^{-1/2},
\]
by \eqref{eq:triangle-marginal} twice. Choose \(\eta>0\) small enough. The final triangle escape therefore gives averaged mass \(\gtrsim R^{-3/4}\) of these three-exit paths with \(q\ge\delta R\) and \(X\ge\eta R\).

Partition \([R,2R]\) into a fixed number of windows shorter than
\(\min(\delta,\eta)R/2\), and retain one window \(J\) with the same lower average. On the fixed measure \(\mu\) of all half-plane returns from \(a\), weighted by \(\kappa^{\ell}\), let \(E_t\) be the deep three-exit event just constructed and put
\[
 Y(\gamma)=\mathop{\mathrm{average}}_{t\in J}{\bf1}_{E_t}(\gamma).
\]
Thus \(\int Y\,d\mu\gtrsim R^{-3/4}\). If \(t<t'\) in \(J\), the simultaneous event has crossings in the order
\[
 (n_1,t),\quad(n_1,t'),\quad(n_2,t),\quad(n_2,t').
\]
Indeed the depth restrictions put \(p_1(t')\) before the \((n_2,t)\) line and \(p_2(t)\) after the \((n_1,t')\) line. Splitting at these unique crossings leaves a prefix, a middle link and a suffix, each of cost \(O(R^{-1/4})\) by triangle escape or the convex diameter bound, and two difference bridges. Therefore
\[
 \mu(E_t\cap E_{t'})\le C R^{-3/4}B_{|t-t'|/d_0}^{2},
 \qquad
 \int Y^2\,d\mu\le C R^{-5/4}.
\]
The convention \(B_0=1\) includes ties. Cauchy--Schwarz now gives
\(\mu\{Y>0\}\gtrsim R^{-1/4}\). All these returns end at a distance between fixed positive multiples of \(R\) from \(a\); monotonicity of \(K_l\) gives \(K_l\gtrsim l^{-5/4}\), completing the proof.

\end{proof}

The following confinement estimate will keep the connectors in the cylinder construction and the later sewing argument inside prescribed channels.

\begin{lemma}[Narrow-channel bridges]\label{lem:narrow-channel-bridges}
Fix an absolute constant $A$ sufficiently large. There are constants $C,R_0>0$ such that, for every strip of height $H\in d_0\mathbb Z_{>0}$ with row-boundary caps normal to one of the directions $n_j$, every starting port on one cap, and every integer $M\ge1$ with $H/M\ge R_0$, the critical mass of bridges whose transverse displacement from the starting port stays at most $4AH/M$ throughout is at least
\[
 e^{-C M} (H/M)^{-1/4}.
\]
In particular, for every fixed transverse tolerance $\tau>0$, restricting the displacement to at most $\tau H$ retains mass comparable to $H^{-1/4}$ for all sufficiently large admissible $H$, with constants allowed to depend on $\tau$.
\end{lemma}
\begin{proof}
Put \(s=H/M\). Choose the \(M-1\) internal cap levels independently in windows of length comparable to \(s\) around the equally spaced levels, so every piece has height between \(s/2\) and \(2s\). At each piece require transverse excursion at most \(A\) times its height and an endpoint shift directed towards the original axis. Reflection retains at least half the mass after the excursion restriction; the diameter tail and the two-sided strip bound make the retained mass at least \(c s^{-1/4}\) when \(A\) is large. Successive starting offsets have absolute value at most \(2As\): each shift points towards zero and has magnitude at most \(2As\). Every point of the concatenation therefore has offset at most \(4As\).

Let \(\mu\) be critical weight on bridges from the prescribed starting port between the two fixed outer caps. For a list \(\mathbf t\) of internal levels let \(E_{\mathbf t}\) be the event that the bridge has the prescribed confined pieces, and define
\[
 Y(\gamma)=\mathop{\mathrm{average}}_{\mathbf t}
           {\bf1}_{E_{\mathbf t}}(\gamma).
\]
Each level is crossed uniquely, so decomposition at a fixed list is injective and
\[
 \int Y\,d\mu\ge c^M s^{-M/4}.
\]
For two lists, split at both levels in each internal window. There are \(M\) intervening pieces of height comparable to \(s\), and \(M-1\) difference bridges within the windows. Dropping the confinement and avoidance constraints bounds the former by strip masses; averaging each latter factor \(B_{|t_i-t_i'|/d_0}\) costs at most \(Cs^{-1/4}\). Consequently
\[
 \int Y^2\,d\mu\le C^M s^{-(2M-1)/4}.
\]
This also covers \(M=1\), when there are no internal levels. Cauchy--Schwarz gives
\(\mu\{Y>0\}\ge e^{-CM}s^{-1/4}\), as claimed. For a fixed tolerance choose a fixed integer \(M\ge4A/\tau\) and combine the lower bound with the strip upper bound.

\end{proof}

\subsection{A barrier estimate on the cylinder}\label{c6:subheading-12}

Take the quotient by \(N\) horizontal lattice steps. We first construct essential polygons in slabs whose height is proportional to the period. Their positive mass will then control the vertical extent of other polygons and arches.

\begin{proposition}[Essential polygons in a slab]\label{prop:slab-essential-polygons}
There are absolute constants $c,C>0$ such that, for every sufficiently large integer $N$ and every horizontal slab $S$ of height $CN$ in the cylinder obtained by quotienting by $N$ horizontal lattice steps,
\[
 \sum_{\substack{\gamma\subset S\\ \gamma\text{ simple essential polygon}}}
       \kappa^{\ell(\gamma)}\ge c.
\]
Here essential means noncontractible in the cylinder, and the polygons are unoriented and unrooted.
\end{proposition}
\begin{proof}
Each endpoint of an arch will first exit a triangle, then cross a narrow tilted strip, and finally exit a second triangle at a common top level. We join these two endpoint collars above that level. A second moment estimate over the collar cuts gives confined arches; joining an upper and a lower arch and averaging over their base level gives the slab bound.

\paragraph{Collars.}

We give the construction with details of aperture restrictions because long-arch confinement in the periodic direction matters here. At a horizontal row boundary of level \(y_0\) in the middle region, take two midpoints \(p,q\) with horizontal separation \(D\in[N/2-\epsilon N,N/2+\epsilon N]\), \(q\) to the right in the lift. We build an arch from \(p\) to \(q\) above the line, of height \(O(N)\) and horizontal span \(<N\). The construction has symmetric independent collars at the endpoints. At \(p\):

\begin{itemize}

\item 

Exit \(T(p,n_0,s_1)\) at \(p_1\) on the side normal to \(n_1\), requiring \(n_0\cdot(p_1-p)\ge\eta N\).

\item 

Make a bridge along \(n_1\) to level \(s_2\) relative to \(p\), ending at \(p_2\), in a narrow tube transversely around the \(n_1\)-line from \(p_1\).

\item 

Exit \(T(p_2,n_1,h-n_0\cdot(p_2-p))\) at \(p_h\) on the side normal to \(n_0\) (level \(y_0+h\)). This size is half-odd in units of \(d_0\) since \(p_2\) is a midpoint in the tilted line.

\end{itemize}

Reflect these instructions at \(q\) to obtain \(q_h\). Figure~\ref{fig:geometry-collars} shows the two collars, their tilted renewal cuts and the common top line.
\begin{figure}[htbp]
\centering
\begin{tikzpicture}[x=1.03cm,y=1.03cm,>=Latex,font=\small,
  region/.style={draw=black!55,thick},
  walk/.style={draw=blue!65!black,very thick},
  cut/.style={draw=black!60,densely dashed}]
  \draw[thick] (-1.1,0)--(11.1,0);
  \draw[cut] (-1.1,3.15)--(11.1,3.15);
  \node[anchor=west] at (-1.05,3.4) {shared level $y_0+h$};
  \node[anchor=south] at (5,.18) {base level $y_0$};

  \path[region,fill=blue!7] (-.63509,0)--(.63509,0)--(0,1.1)--cycle;
  \path[region,fill=orange!13]
    (.18981,.77124)--(.32981,.52876)--(3.57740,2.40376)
    --(3.43740,2.64624)--cycle;
  \path[region,fill=blue!7]
    (3.86824,1.9)--(3.14656,3.15)--(4.58993,3.15)--cycle;
  \draw[cut] (.87203,-.41038)--(-.26519,1.55933);
  \node[anchor=east] at (-.29,1.57) {$s_1$};
  \draw[cut] (4.02702,1.625)--(2.98778,3.425);
  \node[anchor=south east] at (2.98,3.43) {$s_2$};
  \draw[walk] (0,0) .. controls (0,.31) and (.05,.53) .. (.25981,.65)
    .. controls (1.2,1.13) and (2.55,1.87) .. (3.50740,2.525)
    .. controls (3.76,2.66) and (4.1,2.89) .. (4.1,3.15);

  \begin{scope}[xshift=10.3cm,xscale=-1]
    \path[region,fill=blue!7] (-.63509,0)--(.63509,0)--(0,1.1)--cycle;
    \path[region,fill=orange!13]
      (.18981,.77124)--(.32981,.52876)--(3.57740,2.40376)
      --(3.43740,2.64624)--cycle;
    \path[region,fill=blue!7]
      (3.86824,1.9)--(3.14656,3.15)--(4.58993,3.15)--cycle;
    \draw[cut] (.87203,-.41038)--(-.26519,1.55933);
    \draw[cut] (4.02702,1.625)--(2.98778,3.425);
    \draw[walk] (0,0) .. controls (0,.31) and (.05,.53) .. (.25981,.65)
      .. controls (1.2,1.13) and (2.55,1.87) .. (3.50740,2.525)
      .. controls (3.76,2.66) and (4.1,2.89) .. (4.1,3.15);
  \end{scope}

  \draw[very thick,orange!80!black]
    (4.1,3.15) .. controls (4.1,4.1) and (5.9,4.1) .. (5.9,3.15);
  \node at (5,4.17) {connector above the top level};
  \fill (0,0) circle (1.8pt) node[below right=2pt] {$p$};
  \fill (10,0) circle (1.8pt) node[below left=2pt] {$q$};
  \fill (.25981,.65) circle (1.8pt);
  \node at (1,.5) {$p_1$};
  \fill (3.50740,2.525) circle (1.8pt);
  \node at (3.15,1.8) {$p_2$};
  \fill (4.1,3.15) circle (1.8pt) node[above left=2pt] {$p_h$};
  \fill (5.9,3.15) circle (1.8pt) node[above right=2pt] {$q_h$};
  \node[align=center] at (1.5,2.65) {narrow bridge\\along $n_1$};
  \node[align=center] at (8.5,2.65) {reflected\\collar at $q$};
  \draw[<->] (0,-.72)--(10,-.72)
    node[midway,fill=white,inner sep=2pt] {$D=N/2+O(\epsilon N)$};
\end{tikzpicture}
\caption{The two endpoint collars and their completion, shown in the
planar lift. Each collar first exits a small triangle, crosses a narrow
tilted strip between the cuts $s_1,s_2$, and exits a final triangle at
the common height $h$. The tilted cuts are chosen independently at the
two endpoints; only $h$ is shared. The first triangles and the bridges
stay below every allowed top level, and the connector stays above the
chosen top level. These separations recover the pieces at the unique
tilted crossings and the two top crossings. Paths are schematic;
collar widths and vertical clearances are exaggerated.}
\label{fig:geometry-collars}
\end{figure}

Use the following independent windows (except that \(h\) is shared by both collars). Put \(t_*=(1/4-2\epsilon)/d_0\); choose \(s_2/N\) in a small window about \(t_*\), \(h/N\) in \(t_*/2+[0.9,1.1]\delta_h\), \(s_1/N\in[\delta_1,2\delta_1]\), on admissible levels (\(s_1,s_2\) half-odd multiples, \(h\) integer multiple of \(d_0\)). Here \(\epsilon>0\) is small, then \(\delta_h\) sufficiently small compared with \(\epsilon\); the half-width of the \(s_2/N\) window and \(\delta_1\) are both sufficiently small compared with \(\delta_h\). Choose \(\eta>0\) sufficiently small given these, and tube radius a sufficiently small fixed multiple of \(N\) compared with \(\eta N\).

Thus the bridge after the deep exit stays above \(y_0\), and both it and the first triangle stay below the lowest available top level \(y_0+h_{\min}\) by a margin comparable to \(\delta_h N\). Indeed before the last triangle the maximal vertical rise from \(p\) is at most \(s_2/2+O(\delta_1 N+\text{tube radius})\), and \(p_2\)'s height is within this error of \(y_0+s_2/2\). The left collar has horizontal displacements from \(p\) between \(-O(\delta_h N)\) and \((1/4-2\epsilon)N+O(\delta_h N)\), the endpoint near the latter within \(O(\delta_h N)\), by \(n_1=(d_0,1/2)\); similarly on the right. Thus collars are separated, themselves simple by the transversal renewals, and each last triangle still stays strictly above the base (its diameter is \(O(\delta_h N)\)), and terminal horizontal separation is in \([\epsilon N,7\epsilon N]\).

\paragraph{Completion.}
Complete by a half-plane arch joining \(p_h,q_h\) above, with diameter \(\le N/10\). This has no unwanted contacts and the whole arch projects simply.

Averaged over the parameters the product mass \(M_0\) of collars before completion is \(\gtrsim N^{-6/4}\): for each, the averaged deep first exit is \eqref{eq:deep-triangle-exit}, the narrow bridge is supplied by Lemma~\ref{lem:narrow-channel-bridges}, and the last triangle uses \eqref{eq:triangle-escape}. This mass is independent of \(D\).

Without truncating the connector diameter, the sum of product weights including the connector is at least \(c_0(\epsilon N)^{-5/4}M_0\) for each \(D\), with absolute \(c_0>0\), by \(K_l\). The loss by truncation, summed over the integer \(D\)'s, is at most \(C_0 N^{-1/4} M_0\), also with an absolute constant, since for fixed collars varying \(D\) just translates the right endpoint, and we have the completed-path diameter bound. Hence for \(\epsilon\) chosen sufficiently small, at least three quarters of the choices of \(D\) retain averaged mass \(\gtrsim N^{-11/4}\). The concatenation (truncated) is injective for fixed parameters: the top level is crossed precisely twice and on each collar the tilted levels have unique crossings.

\paragraph{Second moment over collar cuts.}
Fix one of the retained separations \(D\) and the endpoints \(p,q\). Let \(\mu_{p,q}\) be critical weight on all simple half-plane arches from \(p\) to \(q\). A cut list is
\(\vartheta=(s_1^p,s_2^p,s_1^q,s_2^q,h)\), with the five independent window laws above. Write \(E_\vartheta\) for the event that an arch admits the prescribed collar decomposition and truncated connector, and set
\[
 Y(\gamma)=\mathop{\mathrm{average}}_\vartheta
            {\bf1}_{E_\vartheta}(\gamma).
\]
The recovery observation gives \(\int Y\,d\mu_{p,q}\gtrsim N^{-11/4}\). To bound its second moment, compare two lists on this same arch measure. Their joint mass is at most
\[
 C N^{-11/4} B_{|h-h'|/d_0}^2
 \prod_{\rm both\ collars} B_{|s_1-s'_1|/d_0}\, B_{|s_2-s'_2|/d_0}.
\]
In fact before the first crossing of the lower of the two top levels, both tilted levels of each list on the \(p\) collar were already crossed uniquely (they are reached strictly below all possible top levels, in the first two pieces, then not crossed again on the collar). Thus chronologically split at these four tilted lines. Prefix to the lowest costs \(O(N^{-1/4})\); from the larger first to the smaller second another \(O(N^{-1/4})\); from the largest to the lower top level, start with normal \(n_1\) and end with normal \(n_0\) within the convex wedge on the proper side of both lines, costing \(O(N^{-1/4})\) by the diameter bound (vertical gap \(\gtrsim N\)). The small intervening bridges yield the indicated difference factors.

Similarly split from \(q\), leaving two vertical bridges of gap \(|h-h'|\) and a top connector bounded by \(K_l=O(N^{-5/4})\) using endpoint separation for the higher-top choice. This proves the bound including ties. Each of the four tilted difference factors averages to \(O(N^{-1/4})\); the shared top-level pair averages to \(O(N^{-1/2})\). Thus
\[
 \int Y^2\,d\mu_{p,q}\le C N^{-17/4},
 \qquad
 \mu_{p,q}\{Y>0\}\ge
 \frac{(\int Y\,d\mu_{p,q})^2}{\int Y^2\,d\mu_{p,q}}
 \gtrsim N^{-5/4}.
\]
This is the confined arch mass on each retained choice of \(D\).

\paragraph{Slab barrier.}
The confined arches can now be joined across a base level and averaged over that level. At least three quarters of the choices of \(D\) are retained for the upper arch. The same holds, after \(D\mapsto N-D\), for the lower arch joining \(q\) to \(p+N\). Their intersection therefore contains a positive fraction of all choices. Joining and projecting, summing over both points with at most bounded overcount, gives mass \(\gtrsim N^{-1/2}\) for polygons crossing the specified \(y_0\) exactly twice at such separated points and staying in our slab.

Let \(\nu\) be critical weight on unoriented unrooted simple essential polygons in the fixed slab. For a base level \(y_0\) in the middle window, let \(F_{y_0}\) be the event just constructed, including its upper and lower arch restrictions. We have \(\nu(F_{y_0})\gtrsim N^{-1/2}\). For two levels \(y_0,y'_0\) the joint mass is \(\lesssim N^{-1/2} B_{|y_0-y'_0|/d_0}^2\): split into upper and lower arches at the higher and lower levels respectively (each costing \(O(N^{-5/4})\) with endpoints given), and two intervening bridges, summing over \(O(N^2)\) endpoint pairs on one level and then bounding bridges separately with other ends free. Each simple arch in the cylinder lifts with endpoint displacement at most \(N\) (larger gives interleaving and intersection with its translate), so has at most two endpoint choices given the start and projected end. For \(Y_0=\mathop{\mathrm{average}}_{y_0}{\bf1}_{F_{y_0}}\), averaging over the window of order \(N\) therefore gives
\[
 \int Y_0\,d\nu\gtrsim N^{-1/2},\qquad
 \int Y_0^2\,d\nu\lesssim N^{-1}.
\]
Cauchy--Schwarz yields \(\nu\{Y_0>0\}\gtrsim1\), proving the slab bound.

\end{proof}

\paragraph{Positivity of the negative polymer polynomial.}
We next obtain a partition-function comparison that turns the slab mass into a penalty for large height. The positivity statement below holds for every integer period, including the smallest ones. Let $D$ be the finite cylinder between horizontal row-boundary caps. Write its height as $H=h d_0$, with $h\ge1$ an integer, and retain distinct edges in the quotient by $N\ge1$ horizontal lattice steps. A simple polygon is an embedded closed curve in this cylinder; it may therefore have two edges when the quotient has parallel edges. Let $\mathcal P$ be the finite set of these polygons. For $0\le\beta'\le3/8$, put
\[
 \alpha=\pi\beta'/3,\qquad x=(2\cos\alpha)^{-1},\qquad
 \omega(L)=x^{|L|}\begin{cases}
 -2\cos(4\alpha),&L\text{ contractible},\\
 -2\cos(2\alpha),&L\text{ essential}.
 \end{cases}
\]
For every $\mathcal A\subseteq\mathcal P$, define
\[
 Q(\mathcal A)=\sum_{\substack{\mathcal F\subseteq\mathcal A\\
                         \text{polygons in }\mathcal F\text{ mutually disjoint}}}
                  \prod_{L\in\mathcal F}\omega(L).
\]
The empty system contributes $1$. We will use positivity not only for the full cylinder but for every set of available polygons.

\begin{auxproposition}[Positivity on all polygon subsets]
For every integer period $N\ge1$, every finite height $h\ge1$, every $0\le\beta'\le3/8$, and every $\mathcal A\subseteq\mathcal P$, the polynomial above satisfies $Q(\mathcal A)>0$.
\end{auxproposition}
\begin{proof}

We first describe the quotient at the short periods. After a lattice translation put the bottom cap at level zero. Set $a=(1,0)$ and $b=(1/2,d_0)$, and denote the centers of the triangles
\[
 \begin{split}
 U_{i,j}&=[ia+jb,(i+1)a+jb,ia+(j+1)b],\\
 D_{i,j}&=[(i+1)a+jb,ia+(j+1)b,(i+1)a+(j+1)b]
 \end{split}
\]
by the same letters. Here $i$ is taken modulo $N$ and $0\le j<h$. The links inside row $j$ are $U_{i,j}D_{i,j}$ and $D_{i,j}U_{i+1,j}$; the links between rows are $D_{i,j}U_{i,j+1}$ for $j<h-1$.

If $N=1$, the two links in each row are distinct parallel edges between its two centers. Each link between rows is a bridge of the finite graph, so no polygon uses it. Consequently the only polygons are the $h$ mutually disjoint essential two-edge row cycles. For any subset $\mathcal A$,
\[
 Q(\mathcal A)=\bigl(1-2\cos(2\alpha)x^2\bigr)^{|\mathcal A|}
             =\left(\frac1{2\cos^2\alpha}\right)^{|\mathcal A|}>0.
\]
This proves the assertion at period one, including both parameter endpoints.

If $h=1$ at any period $N$, the graph is just one essential $2N$-edge row cycle. The only nonempty polymer subset therefore has polynomial
\[
 1-2\cos(2\alpha)x^{2N}\ge 1-2\cos(2\alpha)x^2
 =\frac1{2\cos^2\alpha}>0,
\]
since $0<x\le\kappa<1$ and $\cos(2\alpha)>0$.

Suppose henceforth that $N,h\ge2$. The elementary polygon about the interior lattice point $ia+jb$, $1\le j<h$, has successive centers
\[
 U_{i,j},\ D_{i-1,j},\ U_{i-1,j},\ D_{i-1,j-1},\ U_{i,j-1},\ D_{i,j-1}.
\]
These six centers are distinct modulo $N$, even when $N=2$. The polygon and its enclosed disk have horizontal width $1<N$, so their projection is embedded and contractible. Elementary polygons about adjacent points on an interior horizontal row share centers, and those about $ia+jb$ and $ia+(j+1)b$ share centers when both exist. They therefore form a connected overlap family. Every triangle in the cylinder has a corner on an interior row because $h\ge2$, so every polygon meets at least one member of this family. Thus the incompatibility graph of all polygons, in which two polygons are adjacent when they share a center, is connected for every $N\ge2$.

We recall the flux identity with detached polygons. Start at a bottom boundary port $p$ and assign a path the phase $e^{i\beta' W}$, where $W$ is its total turn. At a previously unvisited center the two outgoing turns are $\pm\pi/3$, and $2x\cos\alpha=1$. If an incoming prefix returns to a center previously reached by a stem, pair the two orientations of the unfinished loop there. They cancel the term stopped at the stem with the corresponding detached polygon occupying that center. The incremental turns from the stem are $\pm4\pi/3$ for a contractible loop, whose source is outside, and $\pm2\pi/3$ for an essential loop. To see the latter value, an essential simple loop has total turn zero in flat cylinder coordinates: under the exponential map to the punctured plane its Jordan-curve turn $\pm2\pi$ is supplied by the winding of the frame. In either case subtract twice the closing turn at the center from the loop's total turn, since departure from the stem turns oppositely. This gives exactly the two cosine factors in $\omega(L)$.

Summing these cancellations gives the signed exit identity
\[
 Q(\mathcal P)=\sum_{\gamma:\,p\to\partial D}
 x^{|\gamma|}e^{i\beta'W(\gamma)}
 Q\bigl(\{L\in\mathcal P:L\cap\gamma=\varnothing\}\bigr).
\]
Here paths stop at their first boundary exit, and length counts visited centers. The real coefficients of all path terms are positive: lifted bridges have turn $0$, and lifted same-cap arches have turn $\pm\pi$. At $\beta'=0$ and $0\le x\le1/2$, the same calculation includes additional defect terms on the right, each with factor $1-2x$ times an incoming prefix weight and the polynomial of its available detached polygons. These terms are nonnegative whenever all subset polynomials are nonnegative.

Continue first from $x=0$ to $x=1/2$ at $\beta'=0$, and then along the critical parameters up to but excluding $\beta'=3/8$. If some subset polynomial first became zero, all subset polynomials would be nonnegative at that parameter. Choose an inclusion-minimal zero subset $\mathcal A$. It is nonempty and all its proper subset polynomials are positive. Every individual polygon weight is strictly negative at this first zero. If $L\notin\mathcal A$ conflicts with some member of $\mathcal A$, the deletion identity gives
\[
 Q(\mathcal A\cup\{L\})
 =Q(\mathcal A)+\omega(L)Q\bigl(\mathcal A\setminus\mathcal N(L)\bigr)<0,
\]
where $\mathcal N(L)$ is the set of polygons meeting $L$. Connectedness of the incompatibility graph therefore forces $\mathcal A=\mathcal P$. In the real flux identity all terms are then nonnegative. A bridge obtained by successively visiting $U_{i,j},D_{i,j}$ in every row has positive weight and deletes at least one elementary polygon. Its residual polynomial is a proper-subset polynomial and hence positive, contradicting $Q(\mathcal P)=0$.

At $\beta'=3/8$, all subset polynomials are nonnegative by continuity, contractible polygon weights are zero, and essential polygon weights are $-\sqrt2\kappa^{|L|}$. Adding a polygon can only decrease a subset polynomial, by the same deletion identity. Thus a zero subset would force $Q(\mathcal P)=0$. In the real flux identity the displayed bridge excludes every essential polygon by separation of the two caps. Its residual polynomial is exactly $1$, since the remaining contractible polygons have zero weight. Its strictly positive contribution again contradicts $Q(\mathcal P)=0$. This completes the proof of strict positivity at every period and on the entire closed interval.
\end{proof}

\paragraph{Deletion ratios and the free-energy bound.}
With positivity on all subsets established, we compare the full partition with the partition obtained by deleting polygons that hit a given path or polygon. At the endpoint parameters let \(Z_H(t)\) use essential polygons only, with factor \(-t\), \(0\le t\le\sqrt2\). Positivity for all subsets continues to hold by multiaffinity (average over deletions). If \(H_H(\gamma;t)\) denotes the polynomial restricted to polymers not hitting \(\gamma\), divided by \(Z_H(t)\), then the following bound measures the cost of the deletion:
\begin{equation}\label{eq:cylinder-deletion-ratio}
 H_H(\gamma;t)\ge\exp\left(t\sum_{\ell\ \mathrm{essential},\ \ell\cap\gamma\ne\varnothing}
                 \kappa^{|\ell|}\right).
\end{equation}
Indeed insert the missing polymers one by one; at each insertion the polynomial decreases by \(t\kappa^{|\ell|}\) times a compatible subpartition, at least \(t\kappa^{|\ell|}\) times the current polynomial, by positivity on all subsets. With one starting point at the bottom the real flux identity becomes
\[
 1=B^{\rm cy}_H/Z_H(\sqrt2)+ c\sum_{\gamma\ \mathrm{arch}}\kappa^{|\gamma|} H_H(\gamma;\sqrt2)
\]
where \(B^{\rm cy}_H\) sums cylinder bridges. For \(H\ge N\), the arch part includes ordinary base returns in a small centered triangle of size comparable to \(N\), so \(B^{\rm cy}_H/Z_H\lesssim N^{-1/4}\) by \(S\).

The bridges themselves have mass at least \(e^{-C H/N}N^{-1/4}\) by Lemma~\ref{lem:narrow-channel-bridges} with \(M\) of order \(H/N\) chosen to fit strictly inside the period. Therefore
\begin{equation}\label{eq:cylinder-free-energy}
 -\log Z_H(\sqrt2)\le CH/N,\qquad H\ge N.
\end{equation}

For the bounded set of periods below the threshold needed by the narrow-channel construction, the explicit bridge above gives $Z_H(\sqrt2)\ge\kappa^{2h}$ when $N,h\ge2$. The same inequality for $N=1$ or $h=1$ follows directly from the displayed row-cycle formulas. Thus the same bound $-\log Z_H(\sqrt2)\le CH/N$ holds after enlarging the absolute constant. The endpoint positivity and deletion monotonicity give the same upper bound for $-\log Z_H(t)$ when $0\le t\le\sqrt2$.

Differentiation gives
\begin{equation}\label{eq:cylinder-integrated-ratio}
 \int_1^{\sqrt2}\sum_{\ell\ \mathrm{essential}}\kappa^{|\ell|}H_H(\ell;t)\,dt
 \le C H/N.
\end{equation}
\paragraph{Height tails.}
The bounded periods also have a uniform slab lower bound: each triangular strip contains the essential $2N$-edge row cycle, of weight $\kappa^{2N}$, whose minimum over these finitely many periods is positive. Taking the slab-height constant large enough to contain a full triangular strip therefore extends the barrier lower bound to them. Consequently the following height estimates use constants independent of the period.

The following consequences will be used in the nesting comparison. They also hold in either other lattice-axis orientation.
\begin{auxproposition}[Cylinder barriers and height tails]
On the infinite cylinder of period $N$, the individual critical weights of essential polygons meeting a fixed horizontal level have bounded sum. From a fixed bottom port in a semi-infinite cylinder, arches reaching height $s$ have mass at most $Ce^{-cs/N}$. For $s\ge1$, all polygons meeting a fixed level and with height span in $[s,2s]$ have total mass at most $\operatorname{poly}(N,s)e^{-cs/N}$, with a fixed polynomial bound.
\end{auxproposition}
\begin{proof}
We combine \eqref{eq:cylinder-deletion-ratio} with Proposition~\ref{prop:slab-essential-polygons} to control vertical span. An essential polygon with vertical height span \(s\), or an arch from the bottom reaching up by \(s\), must hit all short essential polygons in each slab completely traversed. Using order \(s/(CN)\) disjoint slabs and the slab mass bound gives \(H_H(\gamma;t)\gtrsim \exp(c' s/N)\) for \(t\ge1\), \(c'>0\).

Consequently the unadorned mass of essential polygons meeting a given horizontal line on the infinite cylinder is \(O(1)\): sum by spans between \(jN,(j+1)N\), each time using a containing cylinder of height \(O((j+1)N)\) and the last integral bound. The unadorned arch mass from a fixed base midpoint reaching \(s\) in a semi-infinite cylinder is \(\lesssim e^{-c' s/N}\) by the arch flux bound and exhaustion.

We explain also the height-tail bound for unmarked polygons. For $N\ge3$ a polygon of large height can be opened into such an arch at the bottom of its lowest occupied triangle row by bounded-cost surgery.

In that row an upward triangle center \(v\) must occur (any downward triangle there needs a tilted connection to such a center), using both tilted neighbors \(w_\pm\). If one such neighbor uses its other tilted connection to another upward center \(u\), cut that segment \(v,w,u\) (remove \(w\)) and attach \(v,u\) to the base line directly.

Otherwise both neighbors continue north and the upward centers across their other tilted connections are unoccupied (being occupied would force both slants there). Remove \(v\) and attach \(w_\pm\) to the base via these free centers.

In either case the surgery yields a simple arch, still reaching within bounded distance of the same height, with bounded multiplicity given the local surgery site.

For $N=2$, every polygon occupying more than one triangular strip uses exactly one of the two upward centers in its lowest strip. Indeed both upward centers would each have to use both tilted links, forcing the polygon to be the four-edge row cycle. Thus both downward centers are used and continue into the next strip, while the other upward center is unoccupied. Replace the two-link segment through the occupied upward center by the two-link segment through the other one. This gives a simple polygon of the same length and exactly the same minimum and maximum heights. The two alternative segments together form the essential four-edge row cycle, so the replacement changes the winding number by $1$ or $-1$ and takes a contractible polygon to an essential one. The operation is an involution, recovered from the lowest strip, and hence is injective. The already established essential-polygon tail therefore bounds these contractible polygons as well. At $N=1$ every polygon is a row digon and has bounded height span. These two cases complete the height-tail argument at the remaining periods.

Thus polygons meeting a given level with span between \(s\) and \(2s\) have mass \(\le\operatorname{poly}(N,s)e^{-c' s/N}\), by summing arch bounds and surgery locations. Rotation gives the same argument in either other lattice-axis orientation.
\end{proof}
These tails allow the nesting comparison in the next subsection to discard separators with excessive vertical or horizontal span.

\subsection{Passing to local nesting and bulk length}\label{c6:subheading-13}

In the plane let \(Z(R)\) be the positive partition function (including empty) of mutually disjoint simple polygons surrounding a given triangular lattice point, within radius \(R\) of it, factor 2 per polygon, activity \(\kappa^{\rm length}\). We claim
\[
 Z(R)=R^{1/12+o(1)}.
\]
Indeed on the even cylinder with antipodal marked points on the same level we computed \(\mathcal Z=N^{1/6+o(1)}\) for factor 2 on separating polygons. It bounds \(Z(c_1 N)^2\) above for sufficiently small \(c_1>0\) (that is, \(\mathcal Z\ge Z(c_1 N)^2\)). Conversely we can disregard essential polygons at bounded multiplicative cost in an upper bound for \(\mathcal Z\), since they must meet the joining segment and the sum of their single-loop weights is bounded. A contractible separator lifts uniquely around a specified preimage of the marked point on its inside, by simplicity (translated lifts have disjoint interiors). Hence its vertical extent straddles the common level. Those with vertical span above \(N^{1+\xi}\), for any fixed \(\xi>0\), are negligible in individual weight sum by the height tail. For the rest, if horizontal span in the lift exceeds \(N^{1+2\xi}\), project the specified lift around its marked point instead onto a cylinder tilted by \(60^\circ\) with integer period length between \(4 N^{1+\xi}\) and \(5 N^{1+\xi}\). It projects injectively even including the disk inside, by its narrow vertical span, and can be recovered from this polygon by the mark; it meets the level there of the mark and its span normal to the new period is of order at least \(N^{1+2\xi}\). The same height tail discards these as well. Disregarding all discarded polygons costs at most the exponential of twice the sum of their individual weights. All others split into the two types by the point inside, lifting to disjoint plane polygons within each type and supported up to radius \(C N^{1+2\xi}\); systems inject accordingly. Thus
\(
 \mathcal Z\lesssim Z(C N^{1+2\xi})^2
\).
Both comparisons, monotonicity and then \(\xi\downarrow0\) give the claim.

\begin{proposition}[Boundary occupancy]\label{prop:occupancy}
Let $D$ be a finite convex tiled domain and let $e$ be an internal honeycomb link crossing a tiling edge $[y,z]$. If $U_D(e)$ is the critical mass of paths using $e$ with ordered free boundary ports, and $Z_D(y),Z_D(z)$ are the positive enclosing-polygon partitions with factor $2$, then
\[U_D(e)\asymp Z_D(y)+Z_D(z).\]
The comparison constants are absolute.\end{proposition}
\begin{proof}
We compare the two enclosing-polygon partitions with a signed exit sum from the cut link, then recover the complete chords using that link.

\paragraph{Cut-link identity.}
Let \(D\) be a finite convex tiled domain, \(e\) an interior link (identified with the midpoint of a tiling edge \([y,z]\)). Cut \(e\), producing two boundary ends at the same midpoint. Let \(Z_e\) count detached systems of pairwise disjoint simple polygons in this cut domain surrounding that midpoint, factor 2 and critical weight as above. For either choice of cut end as source, the total signed exit mass to the \emph{outer} boundary in the cut domain has real part
\[
 S_e= Z_e+\tfrac1{\sqrt2}(M_y+M_z).
\]
Here \(M_y\) sums a polygon using \(e\) (unoriented, weight without factor 2), having \(y\) inside, together with an exterior system of detached polygons as for \(Z_e\), avoiding that polygon; and similarly for \(M_z\).

Indeed transport prefix flux \(\kappa^{|\eta|}\exp(i\beta W(\eta))\) from this source with any disjoint detached polygons of the specified kind allowed. At a free next triangle center, not previously visited and not occupied by a detached polygon, the two continuations conserve flux with these same polygons. Incoming prefixes returning toward a previously visited center pair by reversing the unfinished polygon there, going back to the first visit. Their turn increments beyond the incoming stem before first visit are total turn of the polygon minus twice its closing turn there (the departure from the stem turns oppositely instead). Thus they are \(\pm4\pi/3\) if the stem came from outside (closing turns convex), \(\pm8\pi/3\) if from inside (closing turns concave); indeed the stem is attached along the unique unused ray and cannot intersect the polygon, so its side is determined by whether the polygon surrounds the cut source. The phases sum to zero outside, and inside both equal \(-1\), canceling the term with the prefix stopped at the incoming stem and that same polygon detached instead, of relative weight 2 times the polygon weight. This partitions the blocked arrivals (at previously visited centers or centers occupied by detached polygons), establishing the exit identity with initial value \(Z_e\). Paths to the outer boundary admit no detached polygons by separation. The other exits end at the opposite cut end and close a polygon once with turn \(\pm2\pi\), real phase \(-1/\sqrt2\) and allowed detached systems necessarily exterior to it. This proves the formula. Also the centered partition functions at \(y,z\) in \(D\) are respectively \(Z_D(y)=Z_e+2M_y,\ Z_D(z)=Z_e+2M_z\) by testing crossings of \([y,z]\). Indeed disjoint polygons about \(y\) either all avoid \(e\), hence all surround the midpoint, or just one uses it (the only possible crossing of this segment), and this one must be innermost since it encloses \(y\) but not \(z\).

\paragraph{Recovering complete chords.}
For an ordinary outer source \(a\), in the cut domain \emph{without} detached loops the same contour works with initial value 1 (stem never inside). Subtract the uncut exit identity to see that the signed mass of complete outer paths from \(a\) using \(e\) equals the signed exit to the cut ends. Summing over all boundary sources and reversing these prefixes, the real mass becomes \(2S_e\). All real weights on complete chords are comparable to unsigned ones. Consequently for free boundary endpoints the unsigned mass of complete paths using \(e\) satisfies
\[
 U_D(e)\asymp Z_D(y)+Z_D(z)
\]
with absolute constants.
\end{proof}

The bounded-factor nesting estimate in the next section will turn this identity into the sharp first-length mass used in the exterior sewing argument.
